%% GENERATED by tools/gen_latex.py -- edit content/ch10.py instead.
%% Build:  latexmk -pdf ch10.tex     (or: pdflatex ch10.tex, twice)
\documentclass[aspectratio=169,11pt,t]{beamer}
\usepackage{itbpro}

\coursetitle{Interpreting What ConvNets Learn}
\coursesubtitle{Four ways to open the box --- and the argument that for ConvNets, the box was never as closed as people say.}
\coursekicker{Chapter 10}
\coursesource{Chollet \& Watson, Deep Learning with Python 3e --- chapter 10}
\courseduration{2.5 hours}
\coursepresenter{Prof. Bambang Riyanto Trilaksono $\cdot$ Rahman Indra Kesuma, S.Kom., M.Cs.}{Lead Instructor $\cdot$ Teaching Assistant}
\coursebrand{AI for Professional \textperiodcentered\ ITB \textperiodcentered\ Ch 10}

\title{Interpreting What ConvNets Learn}
\author{Prof. Bambang Riyanto Trilaksono $\cdot$ Rahman Indra Kesuma, S.Kom., M.Cs.}
\date{}

\begin{document}
\covertitle

\framekicker{Learning outcomes}
\begin{frame}[shrink=25]{Session Objectives}
\leadpar{By the end of this session, participants can:}
\begin{isteps}
\item Build a \textbf{multi-output model} that returns intermediate activations, and read what each depth of layer is doing.
\item Explain \textbf{information distillation} and the rising sparsity of activations with depth.
\item Visualise what a filter responds to using \textbf{gradient ascent in input space}, and say how it differs from ordinary training.
\item Produce a \textbf{Grad-CAM heatmap} and use it to explain --- or debug --- a particular classification.
\item Project a model's \textbf{latent space} to two dimensions to find outliers and semantically ambiguous samples.
\item Say which technique answers which question, so the right one is reached for in a real investigation.
\end{isteps}

\vskip 6pt
\vskip 4pt
\reslink{Book}{Chapter 10 --- full text}{https://deeplearningwithpython.io/chapters/chapter10\_interpreting-what-convnets-learn}
\reslink{Notebook}{01\_intermediate\_activations.ipynb}{https://hendrikarisma.my.id/rs/ai-products-course/notebooks/ch10/01\_intermediate\_activations.html}
\reslink{Notebook}{02\_filter\_visualisation.ipynb}{https://hendrikarisma.my.id/rs/ai-products-course/notebooks/ch10/02\_filter\_visualisation.html}
\reslink{Notebook}{03\_grad\_cam.ipynb}{https://hendrikarisma.my.id/rs/ai-products-course/notebooks/ch10/03\_grad\_cam.html}
\reslink{Notebook}{04\_latent\_space.ipynb}{https://hendrikarisma.my.id/rs/ai-products-course/notebooks/ch10/04\_latent\_space.html}
\reslink{Notebook}{All chapter 10 notebooks}{https://hendrikarisma.my.id/rs/ai-products-course/notebooks/index.html\#ch10}
\reslink{Repo}{Author's official notebook}{https://github.com/fchollet/deep-learning-with-python-notebooks/blob/master/chapter10\_interpreting-what-convnets-learn.ipynb}
\vskip 2pt
\end{frame}

\framekicker{Introduction}
\begin{frame}[shrink=25]{The black-box claim, examined}
{\fontsize{9.4}{12.6}\selectfont\color{fg2}It is often said that deep learning models are black boxes: they learn representations that are difficult to extract and present in human-readable form.\par}\vskip 4pt
\begin{iband}[signal]
Partially true for \textbf{certain types} of model. \textbf{Definitely not true for ConvNets} --- their representations are highly amenable to visualisation, largely because they are \hilite{representations of visual concepts}
\end{iband}
{\fontsize{9.4}{12.6}\selectfont\color{fg2}Since 2013 a wide array of techniques has been developed. This chapter covers the four most accessible and useful.\par}\vskip 4pt
\end{frame}

\framekicker{Introduction}
\begin{frame}[shrink=25]{Why this matters beyond curiosity}
{\fontsize{9.4}{12.6}\selectfont\color{fg2}The motivating question the chapter opens with is a clinical one: \textbf{if a model and a human expert disagree, which of you is seeing the truck?}\par}\vskip 4pt
\begin{iband}[signal]
This is especially relevant where deep learning \textbf{complements human expertise} rather than replacing it --- medical imaging being the example the book names. \hilite{An unexplained disagreement is not usable in that setting.}
\end{iband}
\end{frame}

\framekicker{Introduction}
\begin{frame}[shrink=25]{The four techniques, and what each answers}
\begin{center}
\includegraphics[width=0.92\linewidth,height=0.60\textheight,keepaspectratio]{figs/ch10-four.tex.pdf}
\end{center}
\vskip 3pt{\color{ink3}\fontsize{7.4}{9.4}\selectfont Roughly in order of effort. The first uses your own small ConvNet; the next two use a pretrained Xception.\par}
\end{frame}

\coursesection{01}{Visualising intermediate activations}{What each layer does to one particular image.}

\framekicker{Section 10.1}
\begin{frame}[shrink=25]{What we are looking at}
{\fontsize{9.4}{12.6}\selectfont\color{fg2}The \textbf{activation} of a layer is simply what it outputs for a given input. Visualising activations shows how one input is decomposed into the filters the network learned.\par}\vskip 4pt
\begin{iband}[signal]
Feature maps have three dimensions --- width, height, and channels. Because \textbf{each channel encodes a relatively independent feature}, the right way to look at them is \hilite{one channel at a time, as a 2D image}
\end{iband}
\end{frame}

\framekicker{Section 10.1 $\cdot$ listing 10.3}
\begin{frame}[shrink=25]{A model whose outputs are the layers themselves}
{\fontsize{9.4}{12.6}\selectfont\color{fg2}This is the first genuinely multi-output model in the book: one input, nine outputs --- one per layer activation. Chapter 7's Functional API is what makes it a two-line job.\par}\vskip 4pt
\icode{python}{listing 10.3 --- the activation model}{listings/ch10-01.py}
{\fontsize{9.4}{12.6}\selectfont\color{fg2}This works \textbf{only because the model is Functional}. A subclassed model has no graph to reach into --- exactly the trade-off chapter 7 described.\par}\vskip 4pt
\end{frame}

\framekicker{Section 10.1 $\cdot$ listing 10.4--10.5}
\begin{frame}[shrink=25]{Running it, and plotting one channel}
{\fontsize{9.4}{12.6}\selectfont\color{fg2}Calling \inlinecode{predict()} now returns a \textbf{list} of arrays rather than a single one --- one per layer.\par}\vskip 4pt
\icode{python}{listing 10.4--10.5 --- activations}{listings/ch10-02.py}
\ioutput{listings/ch10-03.txt}
\begin{iband}[signal]
That channel turns out to be a \textbf{diagonal edge detector} --- though your own channels will differ, since the specific filters a layer learns are \hilite{not deterministic}
\end{iband}
\end{frame}

\framekicker{Section 10.1}
\begin{frame}[shrink=25]{Three things visible in the full grid}
\begin{isteps}
\item \textbf{The first layer is a collection of edge detectors.} At that stage the activations retain almost all the information in the original picture.
\item \textbf{Higher layers become increasingly abstract} and less visually interpretable --- encoding concepts like \textit{cat ear} and \textit{cat eye}. They carry less about the image's appearance and more about its class.
\item \textbf{Sparsity increases with depth.} In the first layer every filter fires; further up, more and more are blank --- meaning \hilite{that filter's pattern is not present in this image}.
\end{isteps}
\end{frame}

\framekicker{Section 10.1}
\begin{frame}[shrink=25]{A deep network is an information distillation pipeline}
\begin{center}
\includegraphics[width=0.92\linewidth,height=0.60\textheight,keepaspectratio]{figs/ch10-distill.tex.pdf}
\end{center}
\vskip 3pt{\color{ink3}\fontsize{7.4}{9.4}\selectfont Irrelevant information is filtered out; useful information is magnified and refined.\par}
{\fontsize{9.4}{12.6}\selectfont\color{fg2}This is \textbf{a universal characteristic} of the representations deep networks learn, not something particular to this model.\par}\vskip 4pt
\end{frame}

\framekicker{Section 10.1}
\begin{frame}[shrink=25]{The bicycle test}
{\fontsize{9.4}{12.6}\selectfont\color{fg2}The book offers a striking analogy. After looking at a scene for a few seconds you can remember which abstract objects were in it --- bicycle, tree --- but not their specific appearance.\par}\vskip 4pt
\begin{iband}[signal]
Try drawing a generic bicycle from memory right now. You have seen thousands, and you will almost certainly \hilite{not get it remotely right}. The effect is real, and it is the same distillation: your brain abstracted the input and discarded the detail.
\end{iband}
\end{frame}
\note{Actually do this in the room --- 60 seconds, everyone draws a bicycle. It lands the point better than any slide can.}

\framekicker{Section 10.1 $\cdot$ listing 10.6}
\begin{frame}[shrink=25]{Stitching every channel into one grid}
{\fontsize{9.4}{12.6}\selectfont\color{fg2}Plotting one channel is a spot check. To see the whole picture, every channel of every layer is tiled into a single canvas, layer by layer.\par}\vskip 4pt
\icode{python}{listing 10.6 --- the display loop}{listings/ch10-04.py}
{\fontsize{9.4}{12.6}\selectfont\color{fg2}One canvas per layer, sixteen channels to a row. The next slide explains the standardisation in the middle of that loop, which is the part that makes the result readable.\par}\vskip 4pt
\end{frame}

\framekicker{Section 10.1 $\cdot$ listing 10.6}
\begin{frame}[shrink=25]{Why each channel is standardised first}
\begin{iband}[signal]
The per-channel standardisation matters: without it the faint channels are invisible next to the strong ones, and \hilite{the rising sparsity is exactly what you came to see}
\end{iband}
\end{frame}

\coursesection{02}{Visualising ConvNet filters}{Gradient ascent in input space.}

\framekicker{Section 10.2}
\begin{frame}[shrink=25]{The idea: train the image instead of the weights}
\begin{center}
\includegraphics[width=0.92\linewidth,height=0.60\textheight,keepaspectratio]{figs/ch10-descent-ascent.tex.pdf}
\end{center}
\vskip 3pt{\color{ink3}\fontsize{7.4}{9.4}\selectfont Same machinery, opposite target --- and the sign is flipped.\par}
{\fontsize{9.4}{12.6}\selectfont\color{fg2}Start from a blank (noisy) input and apply gradient descent to the \textbf{image}, maximising the response of one chosen filter. The result is the input that filter likes best.\par}\vskip 4pt
\end{frame}

\framekicker{Section 10.2 $\cdot$ listing 10.7--10.8}
\begin{frame}[shrink=25]{Setting up: a pretrained model, and its layer names}
{\fontsize{9.4}{12.6}\selectfont\color{fg2}The same Xception backbone as chapter 8. First, find out what the convolution layers are called, so their outputs can be retrieved.\par}\vskip 4pt
\icode{python}{listing 10.7--10.8 --- model and layer names}{listings/ch10-05.py}
\ioutput{listings/ch10-06.txt}
{\fontsize{9.4}{12.6}\selectfont\color{fg2}The names follow Xception's block structure from chapter 9, which makes it easy to sample \textbf{shallow, middle, and deep} layers for comparison.\par}\vskip 4pt
\end{frame}

\framekicker{Section 10.2}
\begin{frame}[shrink=25]{The loss: one filter's mean activation}
{\fontsize{9.4}{12.6}\selectfont\color{fg2}Build a feature extractor that stops at the chosen layer, then define a loss that is simply the mean activation of one filter --- with the borders trimmed to avoid edge artefacts.\par}\vskip 4pt
\icode{python}{the loss to maximise}{listings/ch10-07.py}
{\fontsize{9.4}{12.6}\selectfont\color{fg2}Note what this loss is a function of: \textbf{the image}. That is the whole trick.\par}\vskip 4pt
\end{frame}

\framekicker{Section 10.2.1 $\cdot$ listing 10.11}
\begin{frame}[shrink=25]{One ascent step, in TensorFlow}
{\fontsize{9.4}{12.6}\selectfont\color{fg2}This is the book's second low-level gradient loop --- the first was chapter 2. Two details differ from a training step.\par}\vskip 4pt
\icode{python}{listing 10.11 --- gradient ascent step}{listings/ch10-08.py}
\begin{iband}[signal]
\inlinecode{tape.watch(image)} is required because only \inlinecode{Variable}s are watched automatically --- exactly the rule chapter 3 gave. And the \inlinecode{+=} is what makes it \hilite{ascent rather than descent}
\end{iband}
\end{frame}

\framekicker{Section 10.2.2 -- 10.2.3}
\begin{frame}[shrink=25]{The same step in PyTorch and JAX}
{\fontsize{9.4}{12.6}\selectfont\color{fg2}Nothing conceptually new --- this is chapter 3's three gradient idioms applied to a different variable.\par}\vskip 4pt
\begin{itable}{>{\raggedright\arraybackslash}p{0.2068\textwidth}>{\raggedright\arraybackslash}p{0.7332\textwidth}}
\thead{Backend} & \thead{How the gradient is obtained} \\ \midrule
\textbf{TensorFlow} & \inlinecode{tape.watch(image)}, then \inlinecode{tape.gradient(loss, image)}. \\
\textbf{PyTorch} & \inlinecode{loss.backward()}, then read \inlinecode{image.grad}. \\
\textbf{JAX} & Transform the loss function with \inlinecode{jax.grad}, differentiating with respect to the image argument. \\
\end{itable}
{\fontsize{9.4}{12.6}\selectfont\color{fg2}Whichever backend you use, \hilite{the image is the thing being differentiated}, and that is the only unusual part.\par}\vskip 4pt
\end{frame}

\framekicker{Section 10.2.4 $\cdot$ listing 10.14}
\begin{frame}[shrink=25]{The loop that produces one filter's pattern}
{\fontsize{9.4}{12.6}\selectfont\color{fg2}Start from uniform noise centred on 0.5 --- Xception expects inputs in [0, 1] --- and take thirty large steps.\par}\vskip 4pt
\icode{python}{listing 10.14 --- generating a pattern}{listings/ch10-09.py}
\begin{iband}[signal]
A learning rate of \textbf{10.0} would be absurd for training weights. Here it is fine, because \hilite{we are not trying to converge to a minimum} --- only to move the image somewhere the filter likes.
\end{iband}
\end{frame}

\framekicker{Section 10.2.4 $\cdot$ listing 10.15}
\begin{frame}[shrink=25]{\ldots{}and turning the result back into an image}
{\fontsize{9.4}{12.6}\selectfont\color{fg2}The output is a float array whose values are not in [0, 255], so it has to be standardised and clipped before it can be displayed.\par}\vskip 4pt
\icode{python}{listing 10.15 --- deprocess\_image}{listings/ch10-10.py}
{\fontsize{9.4}{12.6}\selectfont\color{fg2}The final crop removes the edge artefacts that gradient ascent tends to accumulate at the borders.\par}\vskip 4pt
\end{frame}

\framekicker{Section 10.2.4}
\begin{frame}[shrink=25]{What the filter banks look like, layer by layer}
\begin{itable}{>{\raggedright\arraybackslash}p{0.188\textwidth}>{\raggedright\arraybackslash}p{0.2444\textwidth}>{\raggedright\arraybackslash}p{0.5076\textwidth}}
\thead{Depth} & \thead{Example layer} & \thead{What the filters encode} \\ \midrule
\textbf{Early} & \inlinecode{block2\_sepconv1} & Simple \textbf{directional edges and colours} --- sometimes coloured edges. \\
\textbf{Middle} & \inlinecode{block4\_sepconv1} & Simple \textbf{textures} made from combinations of edges and colours. \\
\textbf{Deep} & \inlinecode{block8\_sepconv1} & Textures found in natural images: \textbf{feathers, eyes, leaves}. \\
\end{itable}
\begin{iband}[signal]
Chollet's analogy: each layer learns a \textbf{filter bank} such that its inputs can be expressed as a combination of those filters --- \hilite{similar to how a Fourier transform decomposes a signal onto cosines}
\end{iband}
\end{frame}

\framekicker{Section 10.2.2 $\cdot$ listing 10.12}
\begin{frame}[shrink=25]{The PyTorch version, written out}
{\fontsize{9.4}{12.6}\selectfont\color{fg2}Worth seeing in full, because the image has to be marked as requiring a gradient --- which is not something you normally do to an input.\par}\vskip 4pt
\icode{python}{listing 10.12 --- ascent step in PyTorch}{listings/ch10-11.py}
\begin{iband}[signal]
\inlinecode{requires\_grad\_(True)} on an \textbf{input tensor} is the PyTorch equivalent of \inlinecode{tape.watch(image)} --- both say \textit{differentiate with respect to this, even though it is not a parameter}.
\end{iband}
\end{frame}

\framekicker{Section 10.2.3 $\cdot$ listing 10.13}
\begin{frame}[shrink=25]{\ldots{}and the JAX version}
{\fontsize{9.4}{12.6}\selectfont\color{fg2}In JAX the loss function is transformed, so the thing being differentiated is simply whichever argument comes first.\par}\vskip 4pt
\icode{python}{listing 10.13 --- ascent step in JAX}{listings/ch10-12.py}
{\fontsize{9.4}{12.6}\selectfont\color{fg2}Because \inlinecode{compute\_loss(image, filter\_index)} takes the image first, \inlinecode{jax.grad} differentiates with respect to it \hilite{without any extra ceremony at all}\par}\vskip 4pt
\end{frame}

\coursesection{03}{Class activation heatmaps}{Which part of this picture made the model say that?}

\framekicker{Section 10.3}
\begin{frame}[shrink=25]{What a CAM is, and what it is for}
{\fontsize{9.4}{12.6}\selectfont\color{fg2}A \textbf{class activation heatmap} is a 2D grid of scores for a specific output class, computed at every location of an input image, saying how important that location is to that class.\par}\vskip 4pt
\begin{minipage}[t]{0.494\textwidth}
\begin{icardaccent}
\cardh{Debugging a decision}
Particularly a \textbf{misclassification} --- the field is called model interpretability.
\end{icardaccent}
\end{minipage}%
\hfill
\begin{minipage}[t]{0.494\textwidth}
\begin{icardaccent}
\cardh{Locating an object}
The heatmap tells you \textbf{where} in the frame the evidence was.
\end{icardaccent}
\end{minipage}%
{\fontsize{9.4}{12.6}\selectfont\color{fg2}For a dogs-vs-cats model you could generate one heatmap for \textit{cat} and another for \textit{dog} over the same picture.\par}\vskip 4pt
\end{frame}

\framekicker{Section 10.3}
\begin{frame}[shrink=25]{Grad-CAM, in one sentence}
\begin{center}
\includegraphics[width=0.92\linewidth,height=0.60\textheight,keepaspectratio]{figs/ch10-gradcam.tex.pdf}
\end{center}
\vskip 3pt{\color{ink3}\fontsize{7.4}{9.4}\selectfont Weight \textit{how strongly each channel fires here} by \textit{how much this channel matters to the class}.\par}
{\fontsize{9.4}{12.6}\selectfont\color{fg2}The result is a spatial map of \textbf{how intensely the input image activates the class}. The implementation follows Selvaraju et al.\par}\vskip 4pt
\end{frame}

\framekicker{Section 10.3 $\cdot$ listing 10.17}
\begin{frame}[shrink=25]{The worked example: two African elephants}
{\fontsize{9.4}{12.6}\selectfont\color{fg2}The book uses a photograph of an adult elephant and a calf, fed to a pretrained Xception classifier.\par}\vskip 4pt
\icode{python}{listing 10.17 --- preparing the image}{listings/ch10-13.py}
{\fontsize{9.4}{12.6}\selectfont\color{fg2}Everything that follows asks the same question of this picture: \textbf{which pixels made it say \textit{African elephant}?}\par}\vskip 4pt
\end{frame}

\framekicker{Section 10.3.1 -- 10.3.3}
\begin{frame}[shrink=25]{Getting the gradient of the top class}
{\fontsize{9.4}{12.6}\selectfont\color{fg2}Two things are needed: the output of the last convolution layer, and the gradient of the winning class score with respect to that output.\par}\vskip 4pt
\icode{python}{the TensorFlow version}{listings/ch10-14.py}
{\fontsize{9.4}{12.6}\selectfont\color{fg2}PyTorch and JAX versions differ only in the gradient idiom, exactly as in the filter visualisation.\par}\vskip 4pt
\end{frame}

\framekicker{Section 10.3.4 $\cdot$ listing 10.23}
\begin{frame}[shrink=25]{Pooling and weighting to get the heatmap}
{\fontsize{9.4}{12.6}\selectfont\color{fg2}Average the gradient over space to get one importance number per channel, scale each channel by it, then average across channels.\par}\vskip 4pt
\icode{python}{listing 10.23 --- building the heatmap}{listings/ch10-15.py}
\begin{iband}[signal]
Three lines of arithmetic and no new machinery --- Grad-CAM is \hilite{a weighted average, not a new model}
\end{iband}
\end{frame}

\framekicker{Section 10.3.4 $\cdot$ listing 10.24--10.25}
\begin{frame}[shrink=25]{Normalising and superimposing it}
{\fontsize{9.4}{12.6}\selectfont\color{fg2}Clip the negatives, scale to [0, 1], recolour with a colormap, and lay it over the original at 40\% opacity.\par}\vskip 4pt
\icode{python}{listing 10.24--10.25 --- display}{listings/ch10-16.py}
{\fontsize{9.4}{12.6}\selectfont\color{fg2}The result is the picture with the evidence highlighted --- the artefact you would actually show a domain expert.\par}\vskip 4pt
\end{frame}

\framekicker{Section 10.3.4}
\begin{frame}[shrink=25]{The two questions it answers}
\begin{minipage}[t]{0.494\textwidth}
\begin{icardaccent}
\cardh{Why did it say African elephant?}
Because \textbf{these pixels} carried the evidence. If they are the wrong pixels, you have found a spurious correlation.
\end{icardaccent}
\end{minipage}%
\hfill
\begin{minipage}[t]{0.494\textwidth}
\begin{icardaccent}
\cardh{Where is the elephant?}
The heatmap localises the object without ever having been trained to do detection.
\end{icardaccent}
\end{minipage}%
\begin{iband}[signal]
The first question is the one that matters in review: a model that is right \textbf{for the wrong reasons} passes every accuracy check and \hilite{fails the first time the background changes}
\end{iband}
\end{frame}

\coursesection{04}{Visualising the latent space}{Which inputs does the model consider similar?}

\framekicker{Section 10.4}
\begin{frame}[shrink=25]{Every model embeds its inputs in a latent space}
{\fontsize{9.4}{12.6}\selectfont\color{fg2}All deep learning models work by embedding inputs in a \textbf{latent space} --- a manifold where points close together represent semantically similar inputs. On that manifold the classes become separable, and \hilite{that separability is what makes classification possible}\par}\vskip 4pt
{\fontsize{9.4}{12.6}\selectfont\color{fg2}This is chapter 5's manifold hypothesis, now made visible.\par}\vskip 4pt
\end{frame}

\framekicker{Section 10.4}
\begin{frame}[shrink=25]{How to look at it}
\begin{center}
\includegraphics[width=0.92\linewidth,height=0.60\textheight,keepaspectratio]{figs/ch10-latent.tex.pdf}
\end{center}
\vskip 3pt{\color{ink3}\fontsize{7.4}{9.4}\selectfont Take the activations of one layer as coordinates, then project them down to two dimensions.\par}
{\fontsize{9.4}{12.6}\selectfont\color{fg2}Any layer can be used --- \textbf{each encodes a different latent space}, and they become progressively more semantically organised with depth. The usual choice is the layer just before the classifier.\par}\vskip 4pt
\end{frame}

\framekicker{Section 10.4}
\begin{frame}[shrink=25]{What it is actually good for}
\begin{minipage}[t]{0.3253\textwidth}
\begin{icardaccent}
\cardh{Finding outliers}
Points sitting far from every cluster are usually \textbf{mislabelled or corrupt} --- the noisy data of chapter 5.
\end{icardaccent}
\end{minipage}%
\hfill
\begin{minipage}[t]{0.3253\textwidth}
\begin{icardaccent}
\cardh{Finding ambiguity}
Points between two clusters are the samples the model finds \textbf{semantically ambiguous} --- chapter 5's ambiguous features.
\end{icardaccent}
\end{minipage}%
\hfill
\begin{minipage}[t]{0.3253\textwidth}
\begin{icardgood}
\cardh{Improving the dataset}
Both findings feed straight back into \textbf{cleaning the training set and refining the evaluation set}.
\end{icardgood}
\end{minipage}%
\begin{iband}[signal]
Which makes this the technique that closes the loop: it is an interpretation tool whose output is \hilite{a concrete list of samples to go and look at}
\end{iband}
\end{frame}

\framekicker{Section 10.4}
\begin{frame}[shrink=25]{Projecting it, in practice}
{\fontsize{9.4}{12.6}\selectfont\color{fg2}The activations are high-dimensional, so a dimensionality reduction step is needed before anything can be plotted.\par}\vskip 4pt
\icode{python}{recording and projecting activations}{listings/ch10-17.py}
\begin{iband}[signal]
Colour the points by their \textbf{label}, not by the prediction. A point whose colour disagrees with its neighbourhood is \hilite{either mislabelled or genuinely hard} --- and both are worth knowing.
\end{iband}
\end{frame}

\framekicker{Practice}
\begin{frame}[shrink=25]{Interpretation is a loop, not a report}
\begin{center}
\includegraphics[width=0.92\linewidth,height=0.60\textheight,keepaspectratio]{figs/ch10-pipeline.tex.pdf}
\end{center}
\vskip 3pt{\color{ink3}\fontsize{7.4}{9.4}\selectfont The output of every technique here is meant to change something.\par}
\begin{iband}[amber]
A heatmap that nobody acts on is decoration. The value is in the third and fourth boxes: \textbf{a finding, and a decision} --- even when the decision is to accept the behaviour and write it down.
\end{iband}
\end{frame}

\framekicker{Practice}
\begin{frame}[shrink=25]{And what these techniques do not give you}
\begin{ibullets}
\item They explain \textbf{this model on this input}. They are not a general statement about the model's behaviour on data it has not seen.
\item A convincing heatmap is \textbf{not proof of a correct decision} --- a model can look at the right region and still be wrong.
\item None of them turn the model into a set of rules you can audit line by line. They make it \hilite{inspectable, not transparent}.
\end{ibullets}
{\fontsize{9.4}{12.6}\selectfont\color{fg2}Which is worth saying plainly when a stakeholder asks whether the model can be "explained": these techniques answer a narrower question than the one usually being asked.\par}\vskip 4pt
\end{frame}

\framekicker{Summary}
\begin{frame}[shrink=25]{Which technique answers which question}
\begin{itable}{>{\raggedright\arraybackslash}p{0.4888\textwidth}>{\raggedright\arraybackslash}p{0.4512\textwidth}}
\thead{Question you are asking} & \thead{Technique} \\ \midrule
\textit{What is each layer doing to my input?} & \textbf{Intermediate activations} (§10.1) \\
\textit{What pattern does this filter respond to?} & \textbf{Gradient ascent in input space} (§10.2) \\
\textit{Which part of the image drove this decision?} & \textbf{Grad-CAM heatmap} (§10.3) \\
\textit{Which samples does my model confuse, and which are junk?} & \textbf{Latent-space projection} (§10.4) \\
\end{itable}
\end{frame}

\framekicker{Summary}
\begin{frame}[shrink=25]{What has to survive this chapter}
\begin{isteps}
\item \textbf{ConvNets are not black boxes.} Their representations are visual, and therefore visualisable.
\item \textbf{A deep network is an information distillation pipeline}: less about appearance, more about class, with rising sparsity at depth.
\item \textbf{Gradient ascent trains the image, not the weights} --- same machinery, opposite target.
\item \textbf{Grad-CAM is a weighted average}, and it answers \textit{why here?} as well as \textit{where?}
\item \textbf{Latent-space projection finds the samples to go and inspect} --- outliers and ambiguous cases.
\end{isteps}
\vskip 4pt
\reslink{NOTEBOOK}{03\_grad\_cam.ipynb}{https://hendrikarisma.my.id/rs/ai-products-course/notebooks/ch10/03\_grad\_cam.html}
\reslink{NEXT}{Chapter 11 --- Image segmentation}{https://hendrikarisma.my.id/rs/ai-products-course/ch11/index.html}
\vskip 2pt
\end{frame}

\end{document}
